%%%%%%%%%%%%%%%%%%%%%%%%%%%%%%%%%%%%%%%%%
% Important note:
% This template requires the cv.cls file to be in the same directory as the
% .tex file. The cv.cls file provides the resume style used for structuring the
% document.
%
%%%%%%%%%%%%%%%%%%%%%%%%%%%%%%%%%%%%%%%%%

%----------------------------------------------------------------------------------------
%	PACKAGES AND OTHER DOCUMENT CONFIGURATIONS
%----------------------------------------------------------------------------------------

\documentclass{cv} % Use the custom resume.cls style

\usepackage[left=0.75in,top=0.6in,right=0.75in,bottom=0.6in]{geometry} % Document margins
\usepackage{hyperref}
\usepackage{enumitem}
\usepackage{multirow}
\usepackage{ctex} % 支持中文

\newcommand{\tab}[1]{\hspace{.2667\textwidth}\rlap{#1}}
\newcommand{\itab}[1]{\hspace{0em}\rlap{#1}}
\newlist{myitemize}{itemize}{3}
\setlist[myitemize,1]{label=$\circ$,leftmargin=2.8ex}
\setlist[myitemize,2]{label=$\bullet$,leftmargin=2.8ex}
\setlist[myitemize,3]{label=$\diamond$, leftmargin=2.2ex}
\newenvironment{mbi}{\begin{myitemize}
        \setlength{\topsep}{0.6ex}\setlength{\itemsep}{0.16ex}\vspace{0.36ex}}
        {\end{myitemize}\vspace{0.16ex}}

        \newcommand{\mei}{\end{myitemize}\vspace{0.6ex}}
\newenvironment{tmbi}{\begin{myitemize}
        \setlength{\topsep}{0.36ex}\setlength{\itemsep}{0ex}}
        {\end{myitemize}\vspace{1ex}}

\name{白文超} % Your name
\address{\textbf{电话}: (+86) 15161179117~\\
\textbf{邮箱}: \href{mailto: wenchao.bai.chn@gmail.com}{wenchao.bai.chn@gmail.com}~\\
\textbf{个人主页}: \href{https://bai-wenchao.github.io/}{https://bai-wenchao.github.io/}}


\begin{document}

%----------------------------------------------------------------------------------------
%	教育背景
%----------------------------------------------------------------------------------------
\begin{rSection}[MITRed]{教育背景}
%--copy and paste this region  if you need more--
{\textbf{东南大学}. (导师：\href{https://jujuhoo.github.io/}{金嘉晖~副教授})} \hfill {2022年9月 - 至今}
\\ \textbf{博士研究生} @ 计算机科学与工程学院
\hfill { \textbf{GPA}: 3.52/4.0；\textbf{排名}: 4/60 (前7\%)}
\\\\
{\textbf{东北大学}}. \hfill {2018年9月 - 2022年6月}
\\ \textbf{工学学士} @ 软件学院
\hfill { \textbf{GPA}: 4.07/5.0；\textbf{排名}: 3/49 (前6\%)}
\end{rSection}

%----------------------------------------------------------------------------------------
%	实习经历
%----------------------------------------------------------------------------------------
\begin{rSection}[MITRed]{实习经历}
%--copy and paste this region  if you need more--
{\bf 深圳市计算机科学研究院}. (合作导师：\href{https://homepages.inf.ed.ac.uk/wenfei/}{樊文飞~教授}) \hfill {2023年7月 - 2026年6月}\\
{\bf 研究实习生} @ 基础研究部
% \\
% {\bf iFlytek}, Anhui, China \hfill {\em May 2022 - Mar. 2022}\\
% {\bf Big Data Engineer Intern.} @ Intelligent City Business Group. \\
% \\
% {\bf 58.com Inc.}, Beijing, China \hfill {\em Mar. 2021 - Jul. 2021}\\
% {\bf NLP Engineer Intern.} @ AI Lab.
\end{rSection}


\begin{rSection}[MITRed]{研究兴趣}
\mbi
    \item \textbf{数据集成}: 数据质量维护, 数据中心AI;
    \item \textbf{数据系统}: 硬件 (GPU) 加速, 并行模型;
    \item \textbf{数据智能}: 语义查询, 数据智能体.
\mei
\end{rSection}

%----------------------------------------------------------------------------------------
%    荣誉奖项
%----------------------------------------------------------------------------------------
\begin{rSection}[MITRed]{荣誉奖项}
\mbi
    \item {\bf 国家奖学金} (30,000元)，中国教育部。 \hfill {2025年10月}
    \item {\bf SIGMOD Student Travel Award} (1,800美元), SIGMOD及其他赞助商。 \hfill {2025年5月}
\mei
\end{rSection}
%----------------------------------------------------------------------------------------
%    学术服务
%----------------------------------------------------------------------------------------
\begin{rSection}[MITRed]{学术服务}
%--copy and paste this region  if you need more--
\renewcommand{\arraystretch}{1.5}
\begin{tabular}{ ll }
{\bf 助教}: & 算法设计与分析 @ 东南大学 (本科生, 2025年春季学期)\\
{\bf 审稿人}: & T-ITS'24, WWW'26 \\
\end{tabular}
\end{rSection}
%----------------------------------------------------------------------------------------
%    代表性论文
%----------------------------------------------------------------------------------------
\begin{rSection}[MITRed]{代表性论文}
\vspace{-2em}
\nocite{*}
\renewcommand{\refname}{}
\bibliographystyle{plain}
\bibliography{pub}
\end{rSection}

\end{document}
